\cvsection{Expériences professionnelles}

\begin{cventries}
  \cventry
    {Stagiaire ingénieur de recherche en IA}
    {INRIA, équipe Mnémosyne}
    {Bordeaux, France}
    {Mars 2026 -- Septembre 2026}
    {
      \begin{cvitems}
        \item {Analyse d’un besoin défini avec mon tuteur et conception d’une solution agentique aidant les chercheurs à trouver des articles scientifiques.}
        \item {Pilotage technique en autonomie, avec le suivi de mes encadrants, de la conception jusqu’à l’obtention d’une application fonctionnelle.}
        \item {Construction de l’agent sans framework d’orchestration afin de maîtriser ses outils, son état, sa mémoire, son contexte et ses enchaînements de tâches.}
        \item {Intégration de l’API OpenAI, conception des outils de l’agent et mise en œuvre d’un pipeline RAG avec embeddings, recherche sémantique et ChromaDB.}
        \item {Développement du backend Python/FastAPI, d’API REST documentées avec OpenAPI, de la persistance MongoDB et d’une interface React/Next.js.}
        \item {Création d’un outil Python de collecte de données et de webhooks transmettant les résultats de tâches planifiées.}
        \item {Conteneurisation de composants avec Docker et mise en place d’un pipeline CI/CD sur GitHub.}
        \item {Évaluation d’OpenClaw dans un environnement Docker isolé avec une attention portée aux risques liés à l’accès autonome aux outils et aux fichiers.}
        \item {Comparaison de Qwen, Gemma et GLM et de plusieurs quantifications selon leur latence, leur TTFT, la qualité des réponses et la taille du contexte.}
        \item {Définition du protocole expérimental, analyse des résultats et contribution à un article scientifique documentant les performances et limites observées.}
      \end{cvitems}
    }

  \cventry
    {Stagiaire de recherche en apprentissage par renforcement}
    {Université Karunya}
    {Coimbatore, Tamil Nadu, Inde}
    {Juin 2025 -- Sept. 2025}
    {
      \begin{cvitems}
        \item {Conception et comparaison d’une planification de trajectoire déterministe et d’une approche fondée sur l’apprentissage par renforcement pour la navigation autonome d’un robot hospitalier.}
        \item {Déploiement et optimisation du logiciel et du modèle d’apprentissage par renforcement sur un Raspberry Pi embarqué.}
      \end{cvitems}
    }
\end{cventries}
